\subsection*{39. Designing a Real Price Experiment}

A credible experiment specifies the unit (person, account, cohort, region), the offer version, eligibility, presentation, pricing range, primary outcome, guardrails, and duration before data collection. For example, present 2 fully specified packages to comparable eligible groups and track paid purchases, not just button clicks. Hold customer acquisition channel and service promise as constant as possible if the purpose is to estimate a price effect.

If identical customers may see different prices, address fairness and regulatory risk in advance. In some markets randomized pricing is acceptable; in others transparent package tests or sequential price schedules are more appropriate. If the latter are used, acknowledge that changing demand, competition, and reputation may confound the result. Never quietly interpret time-varying sales as a pure price response.

\subsection*{40. Statistical Noise, Sample Size, and Power}

A founder who sees 4 of 20 prospects buy at 1 price and 6 of 20 at another should not declare the 2nd statistically superior. The difference could easily arise by chance or buyer selection. Rigorous testing can require hundreds or thousands of eligible observations for modest conversion changes, depending on the baseline and desired confidence. Small businesses often cannot obtain that sample economically.

Use decision-relevant evidence rather than numerical theater. A preliminary test can establish that some buyers will pay and reveal objections. A 2nd cohort can show whether the transaction is repeatable. To estimate a precise elasticity or optimize several package features, you need much more data and stronger design. The founder's objective at 0 to 1 is to rule out disastrous assumptions, not estimate a 7-decimal optimal fee.

\subsection*{41. Bayesian Price Learning and Sequential Decisions}

Prices can be treated as competing hypotheses rather than permanent identities. Start with a plausible range informed by actual alternatives and costs. Observe purchases, refusals, refunds, and delivery burden; update the range and choose the next informative offer. The Bayesian spirit is to change confidence with evidence, not just follow the last anecdote.

Be cautious with optional stopping and noisy dashboards: watching a metric every hour until a favorable pattern appears biases interpretation. Record the plan, decision thresholds, dates, and deviations. A rigorous founder can make an economically sound choice under uncertainty without pretending every choice has been established by a randomized trial.

\subsection*{42. Conjoint, Gabor-Granger, and Van Westendorp Are Instruments, Not Verdicts}

Customer surveys can reveal price perceptions and attribute trade-offs. Van Westendorp explores perceived cheap and expensive thresholds; Gabor-Granger asks purchase likelihood at explicit prices; discrete-choice conjoint estimates choices among bundles with varying characteristics. Each depends on respondents understanding the actual product, realistic alternatives, and honest treatment of hypothetical bias.

A small, unrepresentative sample of people who ``like the idea'' can produce scientific-looking curves that have little connection to conversion. Use these methods to identify candidate packages and starting price ranges. Prefer revealed behavior and repeated delivery when deciding whether to commit scarce capital. A priced offer is a more serious test than a general preference survey, though its selection and conversion process must still be understood.

\subsection*{43. Test Value Propositions and Prices Separately When Possible}

Suppose conversion rises after the founder simultaneously lowers price, replaces the headline, adds 2 mentor reviews, and recruits from a warmer audience. What caused the improvement? The test has demonstrated that the new configuration can sell; it has not identified price elasticity. In an early company, that may be enough to proceed with the whole configuration, but the causal claim must remain narrow.

When learning which price to use, hold scope and proof as stable as feasible. When learning which \emph{package} buyers prefer, vary the service configuration deliberately. The more intertwined the changes, the less you can infer about any 1 mechanism. Document those limits rather than turning all successes into confirmation of your favorite pricing theory.

\subsection*{44. After the Sale, the Price Experiment Continues}

A price that produces high paid conversion may attract customers who do not complete the product, demand disproportionate support, or feel regret. Higher fees may increase scrutiny and service expectations. Lower fees may reduce commitment or improve accessibility; neither effect is universal. Therefore evaluate at least acquisition, price realization, activation, delivery quality, customer outcomes, contribution, refunds, retention where relevant, and referrals.

In 1-time education, renewal is not always the proper metric. The question is whether graduates would make the same purchase knowing what they know now, whether the capability gain is credible, and whether the next comparable cohort can be recruited at sustainable cost. A sale is the beginning of the proof, not the end.
